% SGA 3, Expose VII_B, English reconstruction -- Loop 1.
% Scope: The remainder of Section 5.1.4, completing the proof of
% Proposition 5.1 and its supplement, on authority local pages 87--88;
% stops immediately before Section 5.1.5.
% Combined-reader pages: 600--601. Re-edition printed pages: 589--590.
% Original-edition marginal locator in scope: 556.
% Editorial notes in scope: 150--151.
% Authority: Polo--Gille Exp7B-13oct24.pdf,
% SHA-256 D303690E0195654BBE7C3DAA5D89EA23E24CD709A3575217234AC19AE8836694.
% Jacob Reinhold's CC BY 4.0 Markdown was consulted as comparison lineage.
% The genuine kernel-pair and exact commutative diagrams are exact
% PDF-derived Loop 1 crops.

现在回到 \SGARefLink{sga3:VIIB:5.1:proposition}{命题~5.1} 的假设 \textup{\SGARefLink{sga3:vii:paragraph:5-1:equation:i}{(i)}}。\footnote{%
\label{sga3:VIIB:nde:150}编者扩充了原文以下部分，并将证明 \SGARefLink{sga3:VIIB:5.1:proposition}{命题~5.1} 所不需要的“补充”\(\overline{I_B^n}=B\cap\overline{J^n}\) 移到了末尾。}
设 \(G\) 为拓扑平坦的无穷小形式 \(k\)-群，\(A\) 为其仿射代数，\(I_A\) 为 \(A\) 的增广理想，\(B\) 为 \(A\) 的闭子代数，并设
\[
  J=AI_B,
  \qquad
  I_B=B\cap I_A.
\]
令 \(H\) 为形式子群 \(\operatorname{Spf}(A/J)\)，并令
\[
  \pi:A\longrightarrow A/J
\]
为投影。假设对每个 \(n\geq1\)，
\[
  A/\overline{J^n}
\]
都在 \(k\) 上拓扑平坦，且 \(B\) 包含于如下映射对的核 \(\widetilde B\)：
\begin{center}
  \phantomsection
  \label{sga3:VIIB:5.1.4:quotient-kernel-pair}
  % SGA3-VII-DIAG-128
  % SGA3-VII-DIAG-128
% Authority: Exp7B-13oct24.pdf, local p. 87.
% Crop: figures/Exp7B_localp087_quotient_kernel_pair.png
% Locator: tex/components/81_expose_VIIB_section514_proposition51_completion_en.tex, Loop 1 line 43.
\begin{tikzpicture}[baseline=(lhs.base)]
  \node (lhs) {$A$};
  \node[right=3cm of lhs] (rhs)
    {$A\widehat\otimes_k(A/J).$};
  \draw[-{Straight Barb[length=4pt,width=4pt]}]
    ([yshift=2.5pt]lhs.east) --
    node[above] {$\tau_1$}
    ([yshift=2.5pt]rhs.west);
  \draw[-{Straight Barb[length=4pt,width=4pt]}]
    ([yshift=-2.5pt]lhs.east) --
    node[below] {$(\operatorname{id}_A\widehat\otimes_k\pi)\Delta_A$}
    ([yshift=-2.5pt]rhs.west);
\end{tikzpicture}

\end{center}

置
\[
  (A_n)=\operatorname{Gr}_J(A),
\]
令
\[
  I_0=I_A/J
\]
为 \(A_0=A/J\) 的增广理想，并如上面的 \textup{(2)} 定义 \((A'_n)_{n\in\mathbb N}\)。令 \((B_n)_{n\in\mathbb N}\) 和 \((\widetilde B_n)_{n\in\mathbb N}\) 分别为 \(\mathbf{Alpg}/k\) 中与 \(B\) 和 \(\widetilde B\) 上由 \(A\) 的滤过所诱导的滤过相联系的对象；这些滤过由理想
\[
  B\cap\overline{J^n},
  \qquad\text{分别}\qquad
  \widetilde B\cap\overline{J^n}.
\]
定义。显然，
\[
  B_n\subset\widetilde B_n\subset A'_n
  \qquad\text{对每个 }n,
\]
以及
\[
  \mathscr B=\prod_{n\geq0}B_n
  \subset
  \widetilde{\mathscr B}
  =
  \prod_{n\geq0}\widetilde B_n
\]
都是
\[
  \mathscr A'=\prod_{n\geq0}A'_n.
\]
的子代数。

另一方面，\(J\) 和 \(\overline{J^2}\) 分别为
\[
  I_B\widehat\otimes_kA,
  \qquad\text{分别}\qquad
  \overline{I_B^2}\widehat\otimes_kA.
\]
在 \(A\) 中的像。因此，映射
\[
  (I_B/\overline{I_B^2})\widehat\otimes_kA
  \longrightarrow
  J/\overline{J^2}
  =
  A_1
  \simeq
  A'_1\widehat\otimes_kA_0
\]
为满射，且经由
\[
  (I_B/\overline{I_B^2})\widehat\otimes_kA_0.
\]
分解。由于
\[
  A_0=k\cdot1\oplus I_0,
\]
下面的交换正合图表明 \(A'_1\) 为 \(I_B/\overline{I_B^2}\) 的像：
\begin{center}
  \phantomsection
  \label{sga3:VIIB:5.1.4:degree-one-image-diagram}
  % SGA3-VII-DIAG-129
  % SGA3-VII-DIAG-129
% Authority: Exp7B-13oct24.pdf, local p. 88.
% Crop: figures/Exp7B_localp088_degree_one_image_diagram.png
% Locator: tex/components/81_expose_VIIB_section514_proposition51_completion_en.tex, Loop 1 line 115.
\begin{tikzcd}[row sep=2.7em,column sep=3em]
(I_B/\overline{I_B^2})\widehat\otimes_k A_0 \arrow[d]
  & \bigl((I_B/\overline{I_B^2})\widehat\otimes_k I_0\bigr)
      \oplus(I_B/\overline{I_B^2})
      \arrow[l,"{\sim}"'] \arrow[d] \\
J/\overline{J^2} \arrow[d]
  & (A'_1\widehat\otimes_k I_0)\oplus A'_1
      \arrow[l,"{\sim}"'] \arrow[d] \\
0 & 0
\end{tikzcd}

\end{center}
于是
\[
  B_1=A'_1.
\]
由于 \(\mathscr A'\) 由 \(A'_1\) 生成，可知对每个 \(n\geq1\)，
\[
  A'_n\subset B_n\subset\widetilde B_n\subset A'_n,
\]
因而
\[
  B_n=\widetilde B_n=A'_n.
\]
\footnote{\label{sga3:VIIB:nde:151}用编者注 149 的记号，这意味着 \(G/H\) 沿单位截面的形式完备化（其仿射代数为 \(\prod_nB_n\)）作为形式簇同构于形式 \(k\)-群 \(N\)。}

最后，由于 \(G\) 为无穷小的，
\[
  \mathfrak r(A)=I_A,
\]
因而理想 \(\overline{I_A^n}\) 趋于零\textup{（参见 \SGARefLink{sga3:VIIB:0.1.2}{0.1.2}）}。更进一步，理想 \(\overline{J^n}\) 也趋于零，所以 \(\widetilde B\) 和 \(B\) 上的诱导滤过是分离的。此外，由于 \(B\) 是 \(A\) 的闭子代数，它关于理想 \(B\cap\overline{J^n}\) 定义的拓扑是完备的。因此，由 \textup{[CA]}，\S~V.7，引理~1\textup{（亦见下文 \SGARefLink{sga3:VIIB:5.1.5}{5.1.5}）}，可得
\[
  B=\widetilde B.
\]
这就完成了 \SGARefLink{sga3:VIIB:5.1:proposition}{命题~5.1} 的证明。

另有如下补充。对每个 \(n\)，
\[
  \overline{J^n}=A\overline{I_B^n}
\]
既是
\[
  A\widehat\otimes_B\overline{I_B^n}
\]
也是
\[
  A\widehat\otimes_B(B\cap\overline{J^n}).
\]
在 \(A\) 中的像。由假设，形式子群 \(H\) 的仿射代数 \(A/J\) 在 \(k\) 上拓扑平坦。因此，由 \SGARefLink{sga3:VIIB:2.4:theorem}{定理~2.4}，态射
\[
  G=\operatorname{Spf}(A)
  \longrightarrow
  G/H=\operatorname{Spf}(B)
\]
是满的且拓扑平坦。从而
\[
  A\widehat\otimes_B\overline{I_B^n}
  =
  \overline{J^n}
  =
  A\widehat\otimes_B(B\cap\overline{J^n}),
\]
这蕴含
\[
  \overline{I_B^n}=B\cap\overline{J^n}
  \qquad\text{对每个 }n.
\]
这一点也可由映射
\[
  \overline{I_B^n}/\overline{I_B^{n+1}}
  \longrightarrow
  A'_n
  =
  (B\cap\overline{J^n})/
  (B\cap\overline{J^{n+1}})
\]
的满射性，以及将下文 \SGARefLink{sga3:VIIB:5.1.5}{5.1.5}\textup{\SGARefLink{sga3:vii:paragraph:5-1:equation:ii}{(ii)}} 用于
\[
  B'_n=\overline{I_B^n},
  \qquad
  B_n=B\cap\overline{J^n}.
\]
而得到。

\phantomsection
\label{sga3:VIIB:5.1.4:supplement-index-source-note}
\noindent\textit{底本说明。}\enspace
在上面以 \(n\) 为指标的公式中，法文 PDF 印作 \(A'_i\)。上下文中的分子和分母，以及对 \(B'_n\) 和 \(B_n\) 的应用，都要求使用 \(A'_n\)；这里采用后者。
